\DeclareMathOperator*{\argmax}{arg\,max}
\DeclareMathOperator*{\argmin}{arg\,min}

\newcommand{\E}{\mathbb{E}}
\newcommand{\prob}{\mathbb{P}}

\newcommand{\Var}{\mathrm{Var}}
\newcommand{\R}{\mathbb{R}}
\newcommand{\spaceg}{G^*}
\newcommand{\simplexg}{\ensuremath{\bar{G^*}}}
\newcommand{\localrad}[2]{\ensuremath{R_\ndata(#1, #2)}}
\def\state{\overrightarrow{X}} % this means that the main state variable will always be X. We use capital X to the random variable 
\def\bara{\bar{\alpha}}
\def\ausiliar{\rho}
\def\barg{\bar{g}}
\def\barsigma{\bar{\sigma}}
\def\lstate{\operatorname{x}} % this means that the main state variable will always be x. We use lower x to indicate values that the random variable takes, as for densities for example.
\def\rmd{\mathrm{d}} %
\def\law{\mathcal{L}}
\def\interval{[0,T]}
\def\dt{\rmd \operatorname{t}}
\def\dbrown{\rmd\!\brown}
\def\brown{\operatorname{B}}
\def\rset{\mathbb{R}}
\def\rsetpos{\mathbb{R}_{\geq 0}}
\def\statedim{d}
\def\targetdist{\mu_{*}}
\def\eqsp{\;}
\def\score{\nabla \log \rho}
\def\realint{\int_{\R^d}}


%Useful commands from my papers with Sylvain Eric and others
\newcommand{\kl}[2]{\operatorname{D}_{\operatorname{KL}}\left(#1||#2\right)}
\newcommand{\W}[2]{\operatorname{W}\left(#1||#2\right)}
\newcommand{\dist}[2]{\operatorname{D}\left(#1||#2\right)}
\newcommand{\dens}[2]{#1\left(#2\right)}
\newcommand{\condd}[3]{\dens{#1}{#3 | #2}}
\newcommand{\utility}[1]{\mathrm{U}(#1)}
\newcommand{\mseloss}[2]{\operatorname{MSE}(#1; #2)}
\newcommand{\measop}[1]{f(#1)}
\def\measnoise{\varepsilon_\lmeas}

\newcommand{\fwmarg}[1]{p_{#1}} % distribution of the foward sde at instant #1
\newcommand{\fwmargd}[2]{\dens{p_{#1}}{#2}} % densitiy of the foward sde at instant #1
\newcommand{\fwtrans}[2]{p_{#2 | #1}} %transition kernel for foward SDE

\newcommand{\mudata}{\mu_{\text{data}}}


\newcommand{\PE}[2]{\mathbb{E}_{#1}\left[#2\right]} % Expected value operator
\newcommand{\PP}[1]{\mathbb{P}\left(#1\right)} % Expected value operator
\newcommand{\PV}[2]{\mathbb{V}_{#1}\left[#2\right]} % Variance operator
\newcommand{\CPE}[3]{\PE{#1}{#3 | #2}} % Conditional expected value

\newcommand{\gausspdf}{\varphi}
\newcommand{\rme}{\mathrm{e}}
\newcommand{\ola}{\overleftarrow}
\newcommand{\ora}{\overrightarrow}
\newcommand{\eqdef}{:=}
\newcommand{\param}{\theta}
\newcommand{\paramdim}{d_\theta}
\def\paramsp{\Theta}
\newcommand{\pihat}{\widehat{\pi}_N^{(\beta,\param)}}
\newcommand{\sched}{\beta}
\def \Nc{\mathcal{N}}
\newcommand{\pdata}{\pi_{data}}


\newcommand{\logp}[1]{\log \tilde{p}_{T - #1} (\x{#1})}
\newcommand{\pforward}[1]{\vec{
p}_{#1}}
\newcommand{\backtildex}[1]{\tilde{X}_{#1}}

\newcommand{\trainedsolution}[1]{X^{\param}_{#1}}
\newcommand{\trainedsolutioninit}[2]{X^{\theta,#2}_{#1}}
\newcommand{\trainedscorex}[1]{s_{\param}(T- #1, \back{#1})}
\newcommand{\trainedscorenov}[1]{s_{\param}(T- #1, \back{#1}^x)}
\newcommand{\dl}{\, d\lambda_{\mathbb{R}^d}}
\newcommand{\ptildex}{\tilde{p}_t(x)}
\newcommand{\pforwardx}[1]{\vec{p}_{#1}}

\newcommand{\gam}{\gamma(x)}
\newcommand{\trainedscore}[1]{s_{\param}(T- #1, \trainedsolution{#1})}

\newcommand{\back}[1]{\overleftarrow{X}_{#1}}
\newcommand{\backscoretilde}[1]{\nabla \log \tilde{p}_{T - #1} (\back{#1})}

\newcommand{\backtildedensityW}[1]{\tilde{p}^{\param}_{#1}}
\newcommand{\backscore}[1]{\nabla \log \pforward{T - #1} (\back{#1})}
\newcommand{\backscorenov}[1]{\nabla \log \pforward{T - #1} (\back{#1}^x)}

\newcommand{\backscoretheta}[1]{\nabla \log \pforward{T - #1} (\trainedsolution{#1})}

\newcommand{\scorep}[1]{\nabla \log \pforward{#1}}
\newcommand{\net}[1]{s_{\param}(T-#1)}

\newcommand{\discsolution}[2][\param]{p_{#2}^{#1}}
\newcommand{\schedulerVE}{g}
\newcommand{\information}{\mathcal{I}}
\newcommand{\pinf}{\pi_{\infty}}

\newcommand{\ind}{\mathds{1}}
\newcommand{\Id}{\mathrm{I}}
\newcommand{\scfc}{\mathsf{s}}

\newcommand{\initerror}{\mathsf{E}_{\operatorname{init}}}
\newcommand{\trainerror}{\mathsf{E}_{\operatorname{train}}}
\newcommand{\discerror}{\mathsf{E}_{\operatorname{disc}}}
\newcommand{\rmL}{\mathrm{L}}
\newcommand{\calF}{\mathcal{F}}
\newcommand{\borelians}[1]{\mathcal{B}(#1)}
\def\tensprod{\otimes}
\def\ndata{n}
\makeatletter
\NewDocumentCommand{\distn}{g g}{%
  d_{\ndata}\IfValueTF{#1}{\left( #1 , #2 \right)}{}%
}
\def\maxray{r_{\ndata}}
\NewDocumentCommand{\distnparam}{g g}{%
  \bar{d}_{\ndata}\IfValueTF{#1}{\left( #1 , #2 \right)}{}%
}
\newcommand{\nballs}[3]{N(#1, #2, #3)}
\newcommand{\chunk}[3]{{#1}^{#2:#3}}
\newcommand{\emin}[1]{\widehat{#1}^{\ndata}}
\newcommand{\tmin}[1]{#1^{\star}}
\NewDocumentCommand{\mfun}{o o}{\operatorname{m}\IfValueTF{#1}{(#2; #1)}{}}
\NewDocumentCommand{\zeromfun}{o o}{\operatorname{g}_{\IfValueTF{#1}{#1}{}}\IfValueTF{#2}{\left(#2\right)}{}}
\NewDocumentCommand{\eintmfun}{o o}{\operatorname{M}_{\ndata}\IfValueTF{#2}{\IfValueTF{#1}{(#2; #1)}{(#2)}}{}}
\NewDocumentCommand{\tintmfun}{o}{\operatorname{M}_{\star}\IfValueTF{#1}{(#1)}{}}
\NewDocumentCommand{\empmeas}{o o}{\mu_{\operatorname{emp}}\IfValueTF{#1}{\IfValueTF{#2}{\left(#1; #2\right)}{\left(#1\right)}}{}}
\NewDocumentCommand{\lpnorm}{m o o}{\left\|#1\right\|_{\operatorname{L}^{#2}(#3)}}
\NewDocumentCommand{\onorm}{m o}{\left\|#1\right\|_{\psi_#2}}
\NewDocumentCommand{\flow}{o o}{\operatorname{T}\IfValueTF{#1}{_{#1}}{}\IfValueTF{#2}{(#2)}{}}
\NewDocumentCommand{\jflow}{o o}{\operatorname{J}_{\operatorname{T}\IfValueTF{#1}{_{#1}}{}}\IfValueTF{#2}{(#2)}{}}
\NewDocumentCommand{\invflow}{o o}{\operatorname{T}^{-1}\IfValueTF{#1}{_{#1}}{}\IfValueTF{#2}{(#2)}{}}
\NewDocumentCommand{\pset}{m}{\mathcal{P}\left(#1\right)}
\NewDocumentCommand{\mlip}{o}{\operatorname{b}\IfValueTF{#1}{\left(#1\right)}{}}
\newcommand{\logdet}[1]{\operatorname{logdet}\left(#1\right)}
\newcommand{\diameter}[1]{\operatorname{Diam}(#1)}
\def\refnfmeas{u}
\newcommand{\pushforward}[2]{{#1}_{\# #2}}
\newcommand{\worstconcentrationvar}[2]{\operatorname{Z}_{\ndata}(#1, #2)}
\newcommand{\gabriel}[1]{\textcolor{red}{GC: #1}}
\newcommand{\tiz}[1]{\textcolor{red}{TF: #1}}
\newcommand{\cvx}[1]{\overline{#1}}
\newcommand{\lrmc}[1]{\cvx{R}_{\ndata}(#1)}
\newcommand{\bernoulli}[1]{\operatorname{Bern}(#1)}
\def\ctesupconcentration{\overline{\operatorname{C}}}
\newcounter{hypH}
\newenvironment{hypH}{
    \refstepcounter{hypH}
    \begin{itemize}
    \item[{\bf H\arabic{hypH}}]
    }
{\end{itemize}}